\documentclass[a4paper, 11pt]{article}
\usepackage[utf8]{inputenc}
\usepackage[T1]{fontenc}
\usepackage[italian]{babel}
\usepackage{geometry}
\usepackage{graphicx}
\usepackage{tabularx}
\usepackage{fancyhdr}
\usepackage{hyperref}
\usepackage[normalem]{ulem}

% Margini
\geometry{top=3cm, bottom=3cm, left=2.5cm, right=2.5cm}

% Intestazione e piè di pagina
\pagestyle{fancy}
\fancyhf{}
\fancyhead[L]{\small \textit{0x06tudio} -- Lettera di Presentazione}
\fancyhead[R]{\small Anno Accademico 2026/2027}
\fancyfoot[C]{\thepage}
\renewcommand{\headrulewidth}{0.4pt}

\newcommand{\gruppo}{0x06tudio}
\newcommand{\codicecapitolato}{C9}
\newcommand{\nomecapitolato}{SpecCheck}
\newcommand{\azienda}{Zucchetti}
\newcommand{\costo}{11.540\,\texteuro}
\newcommand{\scadenza}{30 Aprile 2027}

% Link
\newcommand{\linksott}[2]{\href{#1}{\uline{#2}}}
\newcommand{\urlvalutazione}{https://github.com/0x06tudio/docs/blob/5c52c6d0a37181c3c0dc6ac414c8361ffdf0bb08/out/verbali/candidatura/resoconto_capitolati.pdf}
\newcommand{\urlpreventivo}{https://github.com/0x06tudio/docs/blob/5c52c6d0a37181c3c0dc6ac414c8361ffdf0bb08/out/verbali/candidatura/preventivo_costi_e_impegni_orari.pdf}
\newcommand{\urlsito}{https://example.com}
\newcommand{\urlrepo}{https://github.com/0x06tudio/docs}

% Membri del gruppo
\newcommand{\membri}{%
    Alberto & Beltrame & 2111011 \\
    Pietro & Brocadello & 2147994 \\
    Andrea & Cavalletto & 2116418 \\
    Catalin & Chirilov & 2084626 \\
    Nicolò & Colabove & 2084622 \\
    Davide & Maggio & 2145610 \\
}
% ============================================================

\begin{document}
\thispagestyle{plain}

\begin{minipage}{0.6\textwidth}
    \vspace{-2cm}
    {\huge \textbf{0x06tudio}} \\[0.5em]
    {\large Progetto didattico di Ingegneria del Software} \\[0.3em]
    {\normalsize Università degli Studi di Padova}
\end{minipage}
\begin{minipage}{0.4\textwidth}
    \raggedleft
    \vspace{0pt}
    \includegraphics[width=6cm]{../../media/logo.png}
    \vspace{1cm}
    \small
    \href{mailto:0x06tudio@gmail.com}{0x06tudio@gmail.com}
\end{minipage}

\vspace{1cm}
\section*{Registro delle modifiche}
\begin{tabularx}{\textwidth}{|c|c|X|c|}
    \hline
    \textbf{Versione} & \textbf{Data} & \textbf{Descrizione} & \textbf{Autore} \\ \hline
    0.1.0 & 2026/10/27 & Creazione e stesura bozza iniziale & Pietro Brocadello \\ \hline
\end{tabularx}
\vspace{1cm}
\begin{center}
    {\LARGE \textbf{Candidatura Capitolato \codicecapitolato}}
\end{center}

\vspace{0.5cm}

\noindent Destinatari:\\
Prof. Vardanega Tullio,\\
Prof. Cardin Riccardo.

\vspace{0.5cm}

\noindent Egregi Professori Vardanega e Cardin,\\
con il presente documento il team \textbf{``\gruppo''} vuole formalizzare la propria candidatura per la realizzazione del capitolato del progetto \textbf{\codicecapitolato}:

\begin{center}
    \textbf{\nomecapitolato}
\end{center}

\noindent proposto dall'azienda \textbf{\azienda}.\\
Le principali motivazioni che ci hanno portato a scegliere il presente capitolato sono le seguenti:

\begin{itemize}
    \item Prima motivazione: \textbf{Tema all'avanguardia:} Agentic Coding, LLM Reasoning e verifica automatizzata sono argomenti tecnologicamente molto stimolanti. 
    \item Seconda motivazione: \textbf{Sinergia con l'esame SWE:} Possibilità di testare il tool direttamente sulla documentazione (Casi d'Uso, Requisiti) prodotta dal gruppo per il progetto stesso (ovviamente questa possibilità sara disponibile verso la fine del progetto).
    \item Terza motivazione: \textbf{Riscontro positivo con l'azienda:} che si è mostrata disponibile ad affiancarci durante lo svolgimento del progetto e ci ha già consigliato strumenti ed idee per affrontare il progetto.
\end{itemize}

\noindent Per maggiori dettagli si rimanda al documento di \linksott{\urlvalutazione}{Valutazione dei Capitolati}.

\vspace{0.3cm}

\noindent Come discusso nel documento di \linksott{\urlpreventivo}{Preventivo Costi} abbiamo preventivato un costo complessivo per la realizzazione del progetto di \textbf{\costo}, fissando come termine ultimo per il completamento il \textbf{\scadenza}.

\vspace{0.3cm}

\noindent Tutta la documentazione redatta dal gruppo è pubblicamente disponibile nel seguente \linksott{\urlsito}{sito ufficiale} e nella \linksott{\urlrepo}{repository} dedicata.

\vspace{0.5cm}
\noindent \textbf{\gruppo} è composto dai seguenti membri:

\begin{center}
    \begin{tabular}{lll}
        \textbf{Nome} & \textbf{Cognome} & \textbf{Matricola} \\
        \membri
    \end{tabular}
\end{center}

\vspace{0.5cm}

\noindent Cordiali Saluti,\\
\gruppo.

\end{document}